\scanpage{52}
Je veux revenir sur la condition b), qui peut se préciser beaucoup, si on veut utiliser à fond l'information qu'on avait dans le numéro précédent.

$b_1)$ Pour $x\in\Sigma\deleted{\unsure{(0)}}=\Sigma^{\unsure{*|<}}(0)$, les $z_\alpha\in\Sigma^{<}(1)$ au dessus de $x$ (pour l'application \uline{but} : $\Sigma^{<}(1)\overset{\uline{b}}{\unsure{-|\longrightarrow}}\Sigma^{<}(0)$) sont en correspondance biunivoque avec les germes \inserted{$B_\alpha$} en $x$ de sous-espaces bordants de $\Sigma^{<}(0)$, \deleted{\unsure{avec}} images par $\uline{b}$ des germes de $\Sigma^{<}(1)$ en \deleted{\unsure{les}} $z_\alpha$.

$b_2)$ \deleted{Il existe un} Si $n=\operatorname{card}(\{z_\alpha\})$, il existe un élément \deleted{\unsure{d}} unique de $\Sigma^{<}(n)$, soit $\widetilde{x}$, tel que les \deleted{$\unsure{B_\alpha|z_\alpha}$} $z_\alpha$ soient les images de $\widetilde{x}$ par \inserted{les $n$ applications} de $\widetilde\Sigma(n)$ dans $\widetilde\Sigma(1)$.
\margin{\diagram{\begin{tikzcd}[column sep=large]
\unsure{\widetilde\Sigma^{<}(n)}\arrow[r,phantom,"{\substack{\cdot\\\cdot}}"{description,pos=0.26,yshift=-2.5ex}]\arrow[r,shift left=2ex,"\varphi_{0,n}"]\arrow[r,"\varphi_{1,n}" description]\arrow[r,shift right=2ex,"\varphi_{n-1,n}"']&\widetilde\Sigma^{*}(1)\arrow[d,phantom,"{\rotatebox{90}{$\in$}}" description]\\
&z_\alpha
\end{tikzcd}}}
(Il en résulte qu'il y a une \uline{structure d'ordre} \inserted{\uline{total}} \uline{naturelle} sur l'ens des ``strates de niveau 1'' $B_\alpha$ \inserted{dans $\Sigma^{<}(0)=\Sigma$} passant par $x$)

\deleted{\unsure{$b_2)$}} $b_3)$ Pour tout \deleted{$1\leq d\leq n$, \unsure{les}} $1\leq d\leq n$, considérons l'application
\[
\Sigma^{<}(d)\underset{\substack{\inserted{\text{imm}}\\\deleted{\unsure{\text{vigueur}}}}}{\longrightarrow}\underset{\substack{\rotatebox{90}{$\in$}\\x}}{\Sigma^{<}(0)}
\]
associée à l'inclusion \deleted{\unsure{de de}} $\{d\}\subset\{0,1,--,d\}$, les pts \inserted{$z$} de $\Sigma^{<}(d)$ au dessus de $x$ sont en correspondance biunivoque avec l'ens des parties \inserted{de card $d$} de un \deleted{$\unsure{\Sigma}$} $I$ des germes $B_\alpha$, ou encore l'ens des parties \inserted{de card $d$} de l'ens $I$ des $z_\alpha$, ou encore avec $\mathcal P_d([0,--,n-1])$, en associant : à telle partie $A$ l'application \inserted{str} croissante unique de $\Delta_d=\{0,--d\}$ en

% Lighter-hand foliation at the lower-left corner, not authorial numbering.
\margin{\handwritten{51}}
